\section{Uniform thinning of the actual centered character sums}
\label{sec:centered-character-sums}

\status{CC1--CC20 reconstruct the Dirichlet analytic inputs and the
uniform tail argument behind \texttt{QH.thin} and \texttt{ThinH}.
The real-character zero exclusion is an explicit hypothesis of the
uniform theorem. For a fixed bounded conductor it follows from QW15;
for a host character involving either selected prime it is supplied
by the original \texttt{RealBox} condition. No assertion that all
such boxes exist is made. The simultaneous-pair zero-density and
relative progression estimates remain separate.}

Put $\delta_\chi=1$ for a principal character and $0$ otherwise,
and write
\[
 L_\chi(s)=L(s,\chi),\quad
 F_\chi(s)=(s-1)^{\delta_\chi}L_\chi(s),\quad
 \psi_0(\chi,u)=\sum_{1\le n\le\lfloor u\rfloor}\Lambda(n)\chi(n)
                                      -\delta_\chi u .       \tag{CC.1}
\]
Characters have their usual zero extension at nonunits. Unless stated
otherwise they are primitive of conductor $N\ge1$. A primitive
principal character has $N=1$. All zero sums count analytic
multiplicity. Constants called absolute are independent of the
modulus, character, height and summation endpoint.

\paragraph{CC1: the primitive functional equation and its normalization.}
Let $a\in\{0,1\}$ satisfy $\chi(-1)=(-1)^a$ and put
\[
 \tau(\chi)=\sum_{r\bmod N}\chi(r)e^{2\pi i r/N},\qquad
 \Theta_\chi(t)=\sum_{n\in\mathbb Z}\chi(n)n^a e^{-\pi n^2t/N}.
\]
For primitive nonprincipal $\chi$ the finite Fourier transform is
$\sum_r\chi(r)e^{2\pi i kr/N}=\overline{\chi(k)}\tau(\chi)$.
For a unit $k$ this follows by substitution. For a nonunit $k$,
primitivity supplies a unit $b$ in the kernel of reduction modulo
$N/(k,N)$ with $\chi(b)\ne1$; otherwise the character factors
through that smaller modulus. Multiplication by $b$ leaves the
additive character fixed and forces the Fourier coefficient to vanish.
The unit-reduction map used here is onto by the Chinese remainder
theorem. Parseval applied to this finite transform gives
$|\tau(\chi)|^2=N$.

Periodize the Schwartz function $x^a e^{-\pi t x^2/N}$ on $N\mathbb Z$.
Its Fourier series and all differentiated series converge absolutely.
The Gaussian Fourier transform, and its first derivative when $a=1$,
therefore give the exact identity
\[
 \Theta_\chi(t)=\epsilon_\chi t^{-a-1/2}\Theta_{\bar\chi}(1/t),
 \qquad \epsilon_\chi=\frac{\tau(\chi)}{i^a\sqrt N},
 \qquad |\epsilon_\chi|=1.                              \tag{CC.2}
\]
For $\Re s>1$, termwise Mellin integration gives
\[
 \mathcal L_\chi(s):=
 (N/\pi)^{(s+a)/2}\Gamma((s+a)/2)L_\chi(s)
 =\frac12\int_0^\infty\Theta_\chi(t)t^{(s+a)/2}\frac{dt}{t}.
                                                               \tag{CC.3}
\]
The theta series decays exponentially at infinity; (CC.2) gives
the corresponding decay at zero. Splitting the integral at one
therefore continues it to an entire function and proves
$\mathcal L_\chi(s)=\epsilon_\chi\mathcal L_{\bar\chi}(1-s)$.
For the primitive principal character use the ordinary theta series
and separate its constant term. The two elementary integrals give
the simple poles at $0,1$ and the usual functional equation for
$\pi^{-s/2}\Gamma(s/2)\zeta(s)$.
Thus the continuation and reflection used below follow from the
actual Gaussian and Gauss-sum normalizations.\footnote{These
normalizations agree with NIST DLMF, Section 25.15,
\url{https://dlmf.nist.gov/25.15}. The argument here supplies the
Poisson--Mellin derivation used in this section.}

\paragraph{CC2: polynomial growth on fixed-width strips and a lower bound on the right.}
There are absolute $C,A>0$ such that, for every character modulo
$N$, not necessarily primitive,
\[
 |F_\chi(s)|\le C(N+1)^A(|s|+2)^A
                         \quad(-2\le\Re s\le6),\qquad
 |F_\chi(2+it)|\ge\zeta(2)^{-1}.                       \tag{CC.4}
\]
Here are sufficient details for the growth assertion. For a
nonprincipal character the partial-sum step function $A_0(x)$ is
$N$-periodic and bounded by $N$. Subtract its mean, take a
periodic primitive, subtract the new mean and repeat. After four
integrations by parts, the initial formula
$L(s)=s\int_1^\infty A_0(x)x^{-s-1}\,dx$ is a polynomial in $s$
with bounded boundary coefficients, plus an integral of the form
\[
 s(s+1)\cdots(s+4)\int_1^\infty P_4(x)x^{-s-5}\,dx,
                  \qquad |P_4(x)|\le C_4N^5.
\]
The removed means contribute polynomial terms because the factor
$s+j$ cancels the denominator from the corresponding constant
integral. The remainder converges absolutely for $\Re s>-4$,
and is bounded by $C_4N^5/(\Re s+4)$. This proves (CC.4), after
increasing $A$. The same integration applied to $\lfloor x\rfloor-x$
proves the bound for $(s-1)\zeta(s)$. A principal character modulo
$N$ adds $\prod_{p\mid N}(1-p^{-s})$, whose absolute value on this
strip is at most $\prod_{p\mid N}2p^2\le N^3$ for $N\ge2$.

In $\Re s>1$, the absolutely convergent Euler product gives
\[
 -\frac{L'_\chi}{L_\chi}(s)
       =\sum_{n\ge1}\frac{\Lambda(n)\chi(n)}{n^s},
 \qquad |L_\chi(2+it)|\ge\prod_p(1+p^{-2})^{-1}
                                                   \ge\zeta(2)^{-1}.
                                                               \tag{CC.5}
\]
The principal factor $|1+it|$ is at least one. These bounds also
establish nonvanishing on the open right half-plane.

\paragraph{CC3: local zero counts and logarithmic derivatives.}
Put $\lambda_N=\max(1,\log N)$ and
$\ell_N(t)=\lambda_N+\log(|t|+3)$. Apply Jensen's formula to
$F_\chi$ in the disk of radius four centered at $2+it$, using
(CC.4). The number of zeros in the concentric radius-$13/4$ disk,
with multiplicity, is at most $C\ell_N(t)$.
Removing those zeros by the finite Blaschke construction of QR15
gives, uniformly for $|s-(2+it)|\le5/2$ away from zeros,
\[
 \frac{F'_\chi}{F_\chi}(s)
       =\sum_{|\rho-(2+it)|\le13/4}\frac{m_\rho}{s-\rho}
                                             +O(\ell_N(t)).    \tag{CC.6}
\]
The regular term is bounded by differentiating the logarithm of
the zero-free quotient on the smaller disk. The gaps between the
three radii are fixed, so the implied constant is absolute.
The proof is exactly the rescaled disk argument of QR15; no
global product for the completed function is needed here.

In particular, the number of nontrivial zeros in any height interval
$[t-1,t+1]$ is at most $C\ell_N(t)$: their real parts lie in
$[0,1]$, so their distances from the disk center are at most
$\sqrt5<5/2$. This proves the actual height-radius-one
\texttt{ZerosWindow} estimate, after enlarging the constant to
replace $\ell_N(t)$ by $\log(N(|t|+2))$.
For the near-one arguments we can instead use the center
$5/4+it$ and the radii $7/10,3/5,1/2$ of QR16. The corresponding
growth and Euler-product lower bounds follow from the same
partial-sum formula on $\Re s\ge11/20$. Thus every local input
of QR17--QR18 is available with $\log B$ replaced by a fixed
multiple of $\lambda_N$.

\paragraph{CC4: the boundary line and the real principal character.}
QW14 proves $L_\chi(1)\ne0$ for every nonprincipal character.
For $0<\sigma<1$, the alternating series
$\eta(\sigma)=\sum_{n\ge1}(-1)^{n-1}n^{-\sigma}$ is positive:
pair consecutive terms and then pass to the convergent limit.
Continuation of
$\eta(s)=(1-2^{1-s})\zeta(s)$ gives $\zeta(\sigma)<0$.
The finite principal Euler factors are positive there. Hence
a principal character has no real zero in $(0,1)$.

There are also no zeros at $1+it$, $t\ne0$. On the right of
that line use the coefficientwise nonnegativity
\[
 3+4\cos\theta+\cos2\theta=2(1+\cos\theta)^2\ge0.
                                                               \tag{CC.7}
\]
For primes dividing the modulus all three character coefficients
are zero; on units (CC.7) applies. If $L_\chi$ had a zero of
order $m\ge1$ at $1+it$, the resulting logarithmic-derivative
inequality at $1+\delta$ would have right-hand singular part
at most $(3-4m)/\delta$. A zero of $L_{\chi^2}$ at $1+2it$
only decreases that upper bound; it has no pole there because
$t\ne0$. All remaining terms are bounded for this fixed
modulus and height, by local factorization. Letting $\delta$
decrease to zero is impossible. Together with (CC.5) this puts
every nontrivial zero in the open critical strip.
For a primitive character, (CC.2)--(CC.3) and the absence of
zeros at one also identify the possible zero at zero: it is
simple precisely for nonprincipal even characters. There are
no other zeros in $-1/2\le\Re s\le0$.

\paragraph{CC5: the uniform zero-free region with its exact exception.}
For all characters modulo $N$ there is an absolute $b>0$ such
that a zero in
\[
           \Re s>1-b/\ell_N(\operatorname{Im}s)           \tag{CC.8}
\]
is real, belongs to a real nonprincipal character, and is simple;
across that modulus there is at most one such zero.
Here the analytic inequalities from QR17--QR18 apply without
the arithmetic-family axioms. We spell out their inputs so
that this reuse does not assume a Dirichlet instance of \texttt{LFam}.
Formula (CC.5) supplies the logarithmic derivative. Formula
(CC.6) supplies the local zero expansion and its error.
All nontrivial zeros have real part at most one by CC4,
so omitted zero terms have nonnegative real part to the right
of one. The principal pole contributes exactly $1/(s-1)$.

For a complex character, (CC.7) gives the three-character
inequality used in QR17. For a real character, conjugation
pairs its nonreal zeros and gives the two small-height terms;
at larger heights the principal reciprocal terms are bounded
as in that paragraph. The two numerical contradictions are
the same as (QR.35)--(QR.38), with a new common error constant
$C$. Choosing $b$ smaller than the resulting fixed reciprocal
multiple of $C$ proves the assertion except uniqueness.
For distinct real characters use
$(1+\chi(n))(1+\chi'(n))\ge0$ on units: its four terms
contain only one principal character. Two zeros inside the
chosen width contradict the same near-one estimate. The
one-character two-term inequality excludes multiplicity two.
These are the computations of QR18; their constants depend
only on the local error constant. The real principal case
was excluded separately in CC4. To match the source's closed strip
and denominator $\log(N(|t|+2))$, shrink $b$ by a fixed factor:
$\ell_N(t)$ and this logarithm are comparable uniformly for $N\ge1$,
and the narrower closed strip lies inside (CC.8).

Fix $c_0>0$. If a real nonprincipal $\chi$ has no real zero in
$(1-c_0,1)$, (CC.8) therefore gives, for some $c=c(c_0)>0$,
\[
 0<\beta<1,\quad L_\chi(\beta+i\gamma)=0
 \quad\Longrightarrow\quad
       \beta\le1-\frac{c}{\lambda+\log(|\gamma|+3)}
 \qquad(\lambda\ge1,\ \log N\le\lambda).               \tag{CC.9}
\]
Take, for example, $c\le\min(b,c_0,1)$ after the fixed denominator
comparison. For a nonreal or principal character the exclusion
hypothesis is vacuous. Zeros with $\Re\rho\le0$ satisfy the same
upper bound automatically because $c\le1$ and $\lambda\ge1$;
CC4 excludes zeros with $\Re\rho\ge1$. Thus this supplies the
literal every-zero form of \texttt{pt\_zero\_free}, with the
exceptional-zero hypothesis retained.

\paragraph{CC6: reflect zeros to control small denominators.}
Primitivity is preserved by inversion: factoring $\chi^{-1}$
through a smaller modulus would factor $\chi$ through it.
If $0<\Re\rho<1$ and $L_\chi(\rho)=0$, the gamma factors in
(CC.3) are finite and nonzero there, so
\[
                     L_{\chi^{-1}}(1-\rho)=0.           \tag{CC.10}
\]
The real-character exclusion passes to $\chi^{-1}$, because
a character of square one equals its inverse. Applying (CC.9)
to (CC.10), for $|\Im\rho|\le2$, gives
\[
          \Re\rho\ge\frac{c}{\lambda+\log5}
                         \ge\frac{c}{3\lambda}.         \tag{CC.11}
\]
Thus a zero near zero is controlled by the same exclusion as
a zero near one. This reflection is essential for the
$1/|\rho|$ terms in the explicit formula.

\paragraph{CC7: sum the reciprocal zero weights.}
Let $Z_\chi(T)=\{\rho:0<\Re\rho<1,\ |\Im\rho|<T,
L_\chi(\rho)=0\}$. It is finite by analyticity and CC4.
Put $\kappa=c/3>0$ in (CC.11).
The unit-height count of CC3 then implies
\[
 \sum_{\rho\in Z_\chi(T)}\frac{m_\rho}{|\rho|}
          \le C_2(\lambda+\log T)^2\qquad(T\ge1),       \tag{CC.12}
\]
where $C_2$ depends only on $c_0$.
Indeed the zeros with $|\gamma|\le2$ lie in a fixed number
of unit intervals; their total multiplicity is $O(\lambda)$,
and $1/|\rho|\le\lambda/\kappa$. For $j\le|\gamma|<j+1$,
$j\ge2$, the weight is at most $1/j$ and the count is
$O(\lambda+\log(j+3))$. Summation gives
$O((1+\kappa^{-1})(\lambda+\log T)^2)$, by the elementary
harmonic-sum bounds. One can choose an absolute multiple of
$C(1+\kappa^{-1})$ for $C_2$. This coarser constant is sufficient
for the existential \texttt{psi\_thin} endpoint; no particular
numerical value from the longer source estimate is asserted.

\paragraph{CC8: the literal explicit formula, including the pole at zero.}
There is an absolute $C_7>0$ such that for every primitive
$\chi$, $2\le T\le u$,
\[
 \left|\psi_0(\chi,u)+
       \sum_{\rho\in Z_\chi(T)}m_\rho\frac{u^\rho}{\rho}\right|
 \le C_7\left\{
 \frac{u\log^2(NuT)}T+\log(Nu)
       +\sum_{\rho\in Z_\chi(1)}\frac{m_\rho}{|\rho|}\right\}.
                                                               \tag{CC.13}
\]
We give the contour details behind this source interface.
Apply QR20 to $\Lambda(n)\chi(n)$ on
$\Re s=1+1/\log u$. Its endpoint error is
$O(u\log^2u/T+\log u)$, including integer $u$.

Choose upper and lower heights separately in $[T,T+1]$.
The CC3 count in a fixed-width neighborhood of either height
is $O(\ell_N(T))$. Deleting intervals of radius
$c/\ell_N(T)$ about their ordinates, for a sufficiently small
absolute $c$, removes less than half of $[T,T+1]$.
At a remaining height, (CC.6) bounds $L'/L$ by
$O(\ell_N(T)^2)$ throughout $-1/2\le\Re s\le3$.
Extending the original right-edge integral by at most one unit
at each end costs $O(u\log^2u/T)$, from (CC.5) and
$\sum_{n\ge2}(\log n)n^{-1-1/\log u}=O(\log^2u)$.
Both horizontal edges therefore cost
$O(u\log^2(NuT)/T)$.

On $\Re s=-1/2$, logarithmic differentiation of (CC.3) bounds
$L'/L$ by $O(\log(N(|t|+2)))$: the reflected Dirichlet series
has real part $3/2$, and the two gamma logarithmic derivatives
are $O(\log(|t|+2))$. For completeness that gamma bound follows
from
\[
 \frac{\Gamma'}\Gamma(z)=-\gamma+
       \sum_{j\ge0}\left(\frac1{j+1}-\frac1{j+z}\right).
\]
On a right half-plane, split the series at $2|z|+2$; the initial
part is $O(\log(|z|+2))$ and the tail is $O(1)$.
The recurrence shifts the two fixed real parts here into that
half-plane, a bounded distance from poles. The left edge costs
$O(u^{-1/2}\log^2(N(T+2)))$, absorbed in the displayed error
because $T\le u$.

The principal pole contributes $u$. A nontrivial zero contributes
$-m_\rho u^\rho/\rho$. At zero, write
$-L'/L(s)=-e_0/s+b_0+O(s)$, where $e_0=1$ for a nonprincipal
even character and $0$ otherwise, by CC4. Its residue after
multiplication by $u^s/s$ is $b_0-e_0\log u$.
Formula (CC.6) at zero, after removing the zero there, gives
\[
 |b_0|\le C\log(N+2)+
                  \sum_{\rho\in Z_\chi(1)}m_\rho/|\rho|.
\]
Zeros of larger height in that fixed disk have reciprocal
absolute value at most one and total multiplicity $O(\log(N+2))$;
the possible negative trivial zero in the disk is also a bounded
distance from zero. Thus no unbounded constant at zero is omitted.
Finally replace the good heights by the strict cutoff $|\gamma|<T$.
The added or removed zeros lie in unit intervals near $\pm T$;
each costs at most $m_\rho u/T$, so their total is absorbed.
The residue theorem proves (CC.13), with exactly the source
\texttt{ExplicitFormulaWith} convention.

\paragraph{CC9: choose the height before the conductor.}
Assume the real exclusion in CC5, $\lambda\ge1$,
$\log N\le\lambda$, and $u\ge e$. Put
\[
 v=\log u,\quad s_0=\sqrt v,\quad T=e^{s_0},\quad
 c'=\min(c/3,1),\quad
 E=\exp\left(-\frac{c'v}{\lambda+\sqrt v}\right).
\]
Then $2\le T\le u$, $T^{-1}\le E$, and $uE\ge1$.
For a zero below $T$,
$\lambda+\log(|\gamma|+3)\le\lambda+\sqrt v+2
\le3(\lambda+\sqrt v)$.
Hence (CC.9) gives $u^{\Re\rho}\le uE$.
Using (CC.12), its complete zero sum is bounded by
$C_2uE(\lambda+\sqrt v)^2$.

In (CC.13), $\log(NuT)\le2(\lambda+v)$.
The first error is at most $4uE(\lambda+v)^2$;
the logarithm is at most $uE(\lambda+v)^2$; and the low zero
sum is at most $C_2uE(\lambda+v)^2$.
Consequently
\[
 |\psi_0(\chi,u)|\le
 [C_7(5+C_2)+C_2]u(\log u+\lambda)^2
             \exp\left(-\frac{c'\log u}
                              {\lambda+\sqrt{\log u}}\right).
                                                               \tag{CC.14}
\]
This proves the actual conductor-sensitive estimate by a fixed
height choice. There is no premise of a zero-free critical strip.

\paragraph{CC10: the uniform primitive \texttt{psi\_thin} theorem.}
Collect CC5--CC9. For every $c_0>0$ there are $C,c>0$ such that
for all primitive $\chi$ of conductor $N$, all
$\lambda\ge\max(1,\log N)$, and all $u\ge e$,
\[
 |\psi_0(\chi,u)|\le Cu(\log u+\lambda)^2
                  e^{-c\log u/(\lambda+\sqrt{\log u})},        \tag{CC.15}
\]
provided
\[
 \chi^2=1,\ \chi\ne1\quad\Longrightarrow\quad
     L_\chi(\beta)\ne0\quad(1-c_0<\beta<1).             \tag{CC.16}
\]
The constants are chosen before $N,\chi,\lambda,u$.
This supplies \texttt{QH.psi\_thin}, including its principal
primitive case, and its explicit-formula, zero-window and
zero-free inputs have been established above.

\paragraph{CC11: fixed-character convergence needs no uniform Siegel bound.}
For any fixed modulus, QW15 supplies a positive real-zero
margin for all its real nonprincipal characters. Apply CC10
to the finitely many primitive characters of conductor dividing
that modulus. For each fixed one this gives
\[
                \psi_0(\chi,u)=O_\chi(u/\log^2u).       \tag{CC.17}
\]
Indeed, once $\sqrt{\log u}\ge\lambda$, the exponential
in (CC.15) is at most $e^{-(c/2)\sqrt{\log u}}$, which absorbs
any prescribed power of $\log u$. An imprimitive character
differs from its primitive inducer only at powers of primes
dividing the larger modulus. Their total contribution is
$O_N(\log u)$, so (CC.17) holds for it as well.
The function $\psi_0$ is a locally bounded step function minus
a linear function. Thus
\[
                   \int_1^\infty
                       |\psi_0(\chi,u)|u^{-2}\,du<\infty.
                                                               \tag{CC.18}
\]
This is a fixed-modulus convergence assertion. It makes no
uniform lower-bound assertion as the modulus increases.

\paragraph{CC12: identify the actual nonprincipal centering constant.}
For a nonprincipal character put $a_n=\Lambda(n)\chi(n)$
and $A(u)=\sum_{n\le u}a_n$. Abel summation gives
\[
 \sum_{n\le U}\frac{a_n}{n}
             =\frac{A(U)}U+\int_1^U\frac{A(u)}{u^2}\,du.
                                                               \tag{CC.19}
\]
Equation (CC.17) makes the first term tend to zero and makes
the integral absolutely convergent. For real $\sigma>1$,
the same calculation gives
$-L'/L(\sigma)=\sigma\int_1^\infty A(u)u^{-\sigma-1}\,du$.
Dominated convergence as $\sigma\downarrow1$, justified by
(CC.18), and QW14 nonvanishing identify the limit:
\[
 S_\chi:=\sum_{n\ge1}\frac{\Lambda(n)\chi(n)}n
                         =-L'_\chi(1)/L_\chi(1).        \tag{CC.20}
\]
The series is asserted convergent, not absolutely convergent.
Its value is the original logarithmic-derivative centering.
Subtract (CC.19) at $y$ from its limit and rearrange to obtain,
for every real $y\ge1$,
\[
 G_\chi(y):=\sum_{n\le y}a_n(y/n-1)-yS_\chi
                       =-y\int_y^\infty A(u)u^{-2}\,du.
                                                               \tag{CC.21}
\]
This also verifies the floor convention: if $y$ is an integer,
the newly included summand has coefficient zero.

\paragraph{CC13: the principal centering is exactly $-1-\gamma$.}
Let $R(u)=\sum_{n\le u}\Lambda(n)-u$. From (CC.18) for
the primitive principal character, $R(u)/u^2$ is integrable.
For $\sigma>1$,
\[
 \sigma\int_1^\infty R(u)u^{-\sigma-1}\,du
        =-\frac{\zeta'}\zeta(\sigma)-\frac{\sigma}{\sigma-1}.
\]
The Laurent expansion
$\zeta(s)=1/(s-1)+\gamma+O(s-1)$ gives the limit $-1-\gamma$.
This constant also follows directly from the sum-minus-integral
definition $\gamma=\lim_n(\sum_{j=1}^n1/j-\log n)$ in the
Euler summation formula for $\zeta$. Dominated convergence yields
\[
                      \int_1^\infty R(u)u^{-2}\,du=-1-\gamma.
                                                               \tag{CC.22}
\]
The finite integral is
$\int_1^y R(u)u^{-2}\,du
=\sum_{n\le y}\Lambda(n)(1/n-1/y)-\log y$.
Subtract it from (CC.22) to get
\[
 G_1(y):=\sum_{n\le y}\Lambda(n)(y/n-1)
                    -y(\log y-1-\gamma)
                       =-y\int_y^\infty R(u)u^{-2}\,du.
                                                               \tag{CC.23}
\]
Thus the principal and nonprincipal constants have both been
identified, without borrowing a uniformly quantified
Siegel--Walfisz estimate to justify convergence.

\paragraph{CC14: match the literal centered functions.}
The source $G_0(\chi,y)$ is $G_1(y)$ when $\chi$ is principal
and $G_\chi(y)$ otherwise. For primitive characters (CC.21)
and (CC.23) therefore give exactly
\[
             G_0(\chi,y)=-y\int_y^\infty
                           \psi_0(\chi,u)u^{-2}\,du.    \tag{CC.24}
\]
The integral exists absolutely. Primitivity in the principal
case is essential: it makes $N=1$, so its von Mangoldt sum
has no deleted Euler factors.
The host function \texttt{An.gChi} uses
$\max(y/n-1,0)$ on $1\le n\le\lfloor y\rfloor$.
On this finite interval $n\le y$, so the maximum is $y/n-1$.
Its nonprincipal centering has the sign $+yL'/L(1)$ from
(CC.20), and its principal term is precisely (CC.23).
Hence it equals $G_0$ on these actual arguments. This supplies
\texttt{QH.g\_integral} and the function-identification clause
used by \texttt{thinH}.

\paragraph{CC15: a uniform integral estimate above the exponential cutoff.}
Fix $\epsilon,c,C,A,c'>0$. Suppose
\[
 y\ge\tfrac13e^{x^{\epsilon/3}},\qquad
 1\le\lambda\le Ax^{\epsilon/5}\log^2x,
\]
and a measurable function $\phi$ satisfies the bound on the
right side of (CC.15) for every $u\ge y$.
For all sufficiently large $x$, independently of $\lambda,y,\phi$,
\[
 \left|y\int_y^\infty\phi(u)u^{-2}\,du\right|
                          \le y e^{-c'x^{\epsilon/8}}.   \tag{CC.25}
\]
Here is a short proof with the uniform quantifiers visible.
Put $V=\log y\ge\tfrac12x^{\epsilon/3}$.
Since $61/300>1/5$, the power $x^{\epsilon/300}$ eventually
dominates $\log^2x$ and the fixed constant $A$.
Consequently, for every $v\ge V$,
$\lambda\le v^{61/100}$ and $v\ge1$. Therefore
\[
 (v+\lambda)^2e^{-cv/(\lambda+\sqrt v)}
                    \le4v^2e^{-(c/2)v^{39/100}}.        \tag{CC.26}
\]
The substitution $u=e^v$ in the norm majorant reduces its
integral to the integral of the left side of (CC.26).
Factor half the exponential from the right side at $V$:
\[
 \int_V^\infty4v^2e^{-(c/2)v^{39/100}}\,dv
 \le K_c e^{-(c/4)V^{39/100}},\qquad
 K_c=4\int_0^\infty v^2e^{-(c/4)v^{39/100}}\,dv<\infty.
\]
The last integral is finite by $w=v^{39/100}$, or directly
by exponential decay. The exponent at $V$ is at least a
positive constant times $x^{13\epsilon/100}$.
Since $13/100-1/8=1/200>0$, this absorbs $CK_c$ and any
prescribed $c'x^{\epsilon/8}$. This proves (CC.25).
Measurability and the displayed integrable majorant also prove
the claimed absolute integrability. The source's more general
totalized integral convention assigns zero to a nonintegrable
integrand; that unused case satisfies its numerical inequality
trivially. For $\phi=\psi_0$, measurability was proved in CC11.

\paragraph{CC16: the exact uniform thinning theorem.}
Fix $\epsilon,A,c',c_0>0$. There is $x_0$ such that for every
$x\ge x_0$, every primitive character $\chi$ with
\[
       \log N\le Ax^{\epsilon/5}\log^2x,
\]
satisfying (CC.16), and every
$y\ge\tfrac13e^{x^{\epsilon/3}}$, one has
\[
                        |G_0(\chi,y)|
                                \le y e^{-c'x^{\epsilon/8}}.
                                                               \tag{CC.27}
\]
Choose $\lambda=\max(1,\log N)$. For large $x$ it is at most
$(A+1)x^{\epsilon/5}\log^2x$. Apply CC10 with its constants
chosen only from $c_0$, then CC15 and the exact identity CC14.
Increasing $x_0$ ensures $y\ge e$ and all displayed thresholds.
The resulting threshold depends only on the four fixed
parameters. This is \texttt{QH.thin}, including its universal
quantifier over every $y$ above the cutoff.

\paragraph{CC17: a primitive real character has restricted conductor exponents.}
If a primitive character has square one, every odd prime
occurs in its conductor to exponent at most one, and the
exponent at two is at most three.
For an odd prime $p$, the kernel of
$(\mathbb Z/p^k\mathbb Z)^\times\to(\mathbb Z/p\mathbb Z)^\times$
is a finite $p$-group. A homomorphism from it to $\{\pm1\}$
is trivial, since its image order divides both a power of $p$
and two. At two, every unit congruent to one modulo eight
is a square modulo $2^k$. To prove the latter, lift an odd
square root from modulus $2^j$ to $2^{j+1}$ by replacing
$z$ with $z+2^{j-1}$ if necessary; for $j\ge3$ this flips
exactly the next binary digit of $z^2$.
A square-one character kills these squares. Chinese remainders
then make it factor through $8\operatorname{rad}(N_{\rm odd})$.
Primitivity proves the asserted exponent bounds.

Consequently, if $8\mid M$, $q,d$ are primes, and every prime
dividing $N$ divides $M$ or equals $q$ or $d$, then
\[
                                  N\mid Mqd.             \tag{CC.28}
\]
This is the actual \texttt{real\_cond} assertion, with its
prime-support hypotheses. If $q,d\nmid N$, the stronger
$N\mid M$ follows immediately. No exponent bound is imposed
on a general complex character.

\paragraph{CC18: bounded conductors have one common real-zero exclusion interval.}
For each fixed integer $B_0\ge0$, there is $c_s>0$ such that
every nonprincipal character modulo $N\le B_0$
has no real zero in $(1-c_s,1)$.
For each such character QW14 makes its analytic function
nonzero at one; continuity gives a neighborhood of one free
of zeros. There are finitely many moduli and characters, so
take the minimum of finitely many positive radii. When the
set is empty, take any positive radius. This is the same
finite-minimum argument as QW15, now for all nonprincipal
characters with the precise \texttt{small\_zeros} quantifier;
primitivity and reality are not needed in this bounded-modulus
statement.
Its dependence on $B_0$ is allowed; it is not a uniform
Siegel-type bound for unbounded conductors.

\paragraph{CC19: install the actual host thinning statement.}
Fix $B\in\mathbb R$ and $0<\epsilon<1/52$, as in the
literal \texttt{ThinH} quantifier. Let $H$ be a supplied host,
and write $Q_S$ for its fixed
cyclotomic factor and $q,d$ for its two prime factors.
The actual conductor is $Q_Sqd$. A primitive host character
$\chi$ has $N\mid Q_Sqd$: it is the primitive inducer of
the character associated to its host index. This is the
same conductor compatibility already used in HC9.
Assume
\[
 Q_{S,\rm odd}\le B,\qquad
 \log(Q_Sqd)\le2x^{\epsilon/5}\log^2x,
\]
and the source \texttt{RealBox}: each real primitive host
character involving $q$ or $d$ is nonzero for real
$\beta>1-1/400$.

Apply CC18 with $B_0=8\lfloor B\rfloor$ and put
$c_0=\min(c_s,1/400)>0$.
For a real nonprincipal host character, if $q\mid N$ or
$d\mid N$, the given \texttt{RealBox} supplies (CC.16).
Otherwise CC17 and the prime support of $Q_Sqd$ give
$N\mid8Q_{S,\rm odd}$ and hence $N\le8\lfloor B\rfloor$;
CC18 supplies (CC.16). Nonreal and principal characters
need no such exclusion. Every host character also satisfies
$\log N\le\log(Q_Sqd)$.
Now apply CC16 with $A=2$ and $c'=1$, and identify its
function with \texttt{An.gChi} by CC14. This proves the
literal \texttt{ThinH}, with its outer constant chosen as one.
The threshold depends on $B,\epsilon$, and the statement is
uniform over all supplied hosts, primitive host characters
and $y$ above the original exponential cutoff.
Thus the thinning premise of HC9 is supplied. The construction
of a host satisfying \texttt{RealBox}, simultaneous
\texttt{PairEx}, the relative progression theorem and the
cohomological windows are not inferred from this implication.

\paragraph{CC20: a quantitative derivative consequence for the same centered sums.}
The proof also gives a useful strengthening above the same
exponential cutoff. Under the hypotheses of CC16, for any
fixed $b>0$ and all sufficiently large $x$,
\[
 |G_0(\chi,y_2)-G_0(\chi,y_1)|
       \le (y_2-y_1)e^{-b x^{\epsilon/8}}
 \quad\left(\tfrac13e^{x^{\epsilon/3}}\le y_1\le y_2\right).
                                                               \tag{CC.29}
\]
Away from integers, differentiate (CC.24):
\[
           G'_0(\chi,y)=\psi_0(\chi,y)/y+G_0(\chi,y)/y.
\]
The same identity holds for the right derivative at integers.
Each ramp is continuous there, so $G_0$ is locally absolutely
continuous and piecewise continuously differentiable.
The argument of CC15 applied pointwise to (CC.15) bounds
$|\psi_0(\chi,y)|/y$ by $e^{-2b x^{\epsilon/8}}$, uniformly
over all $y$ above the cutoff, after increasing the threshold.
Use CC16 with $c'=2b$ for the second term. For large $x$,
their sum is at most $e^{-b x^{\epsilon/8}}$.
Integration on $[y_1,y_2]$ proves (CC.29), including intervals
crossing integers. This is an additional human consequence of
the proved bounds; it does not substitute for the separate
square-root estimates on the earlier polynomial-height window.

